\chapter{Meningioma}

\section*{Definition and Overview}
A meningioma is a primary central nervous system (CNS) tumour that arises from the meninges—the protective membranes enveloping the brain and spinal cord. Specifically, these tumours originate from the arachnoid cap cells of the arachnoid villi. Meningiomas are typically slow-growing and represent the most common primary brain tumour diagnosed in adults. Although the vast majority (approximately 80\% to 85\%) are histologically benign and classified as World Health Organization (WHO) Grade I, they can still cause significant clinical symptoms and neurological deficits by exerts pressure on adjacent brain tissue, cranial nerves, and blood vessels. Early identification and classification are critical for determining the optimal balance between active surveillance and therapeutic intervention.

\section*{Epidemiology}
Meningiomas account for approximately 37\% to 40\% of all primary brain and CNS tumours, making them the most prevalent primary intracranial neoplasm. The incidence increases progressively with age, with the peak diagnosis occurring in individuals aged 60 to 80 years. They are exceptionally rare in children, where they are typically associated with genetic syndromes. Meningiomas exhibit a strong female-to-male predilection, occurring two to three times more frequently in women than in men, a discrepancy that is even more pronounced during the reproductive years. In Australia, the age-standardised incidence is estimated at 7 to 9 cases per 100,000 individuals per year. 

\section*{Aetiology and Pathophysiology}
While most meningiomas are sporadic, several environmental and genetic factors have been identified in their pathogenesis.
\begin{itemize}
    \item \textbf{Genetic alterations:} The most common genetic driver is the inactivation of the neurofibromin 2 (\textit{NF2}) tumour suppressor gene on chromosome 22q12, which encodes the Merlin protein. This mutation is found in up to 60\% of sporadic meningiomas and is a hallmark of the hereditary syndrome Neurofibromatosis Type 2. Non-\textit{NF2} mutated tumours often exhibit alterations in other genes, such as \textit{TRAF7}, \textit{AKT1}, \textit{KLF4}, and \textit{SMO}.
    \item \textbf{Ionising radiation:} Exposure to ionising radiation is the most established environmental risk factor. Individuals who received cranial radiotherapy during childhood (e.g., for leukaemia or tinea capitis) have a significantly increased risk of developing meningiomas, often multiple, decades later.
    \item \textbf{Hormonal factors:} The presence of progesterone, oestrogen, and androgen receptors on meningioma cells explains why these tumours often accelerate in growth during pregnancy or with the use of certain exogenous hormones (such as high-dose cyproterone acetate or hormone replacement therapy).
    \item \textbf{Obesity:} Clinical studies indicate that elevated body mass index (BMI) correlates with an increased risk of developing meningiomas, possibly mediated through chronic inflammatory states or altered systemic hormone levels.
\end{itemize}

\section*{Clinical Features}
Meningiomas are frequently asymptomatic in their early stages and are increasingly discovered incidentally during neuroimaging for unrelated symptoms. When symptoms do occur, they are primarily driven by mass effect and location.
\begin{itemize}
    \item \textbf{Headaches:} Typically progressive, worse in the morning or when lying down, and may be accompanied by nausea or vomiting due to increased intracranial pressure.
    \text{Seizures:} New-onset focal or generalised seizures, resulting from local cortical irritation by the tumour.
    \item \textbf{Focal neurological deficits:} Progressive weakness or numbness in the limbs, difficulty with balance and coordination, or cranial nerve deficits such as double vision, facial numbness, or hearing loss.
    \item \textbf{Visual changes:} Blurring, double vision (diplopia), or visual field loss, commonly observed with parasellar or sphenoid wing meningiomas compressing the optic nerve or chiasm.
    \item \textbf{Cognitive and personality changes:} Frontal lobe meningiomas can present with subtle personality shifts, apathy, memory impairment, or executive dysfunction, which are sometimes misdiagnosed as depression or dementia.
    \item \textbf{Anosmia:} Loss of the sense of smell, typical of olfactory groove meningiomas.
\end{itemize}

\section*{Differential Diagnosis}
When evaluating a dural-based mass on imaging, several alternative pathologies must be considered:
\begin{itemize}
    \item \textbf{Dural metastases:} Secondary cancers spreading to the meninges from primary sites such as the breast, lung, kidney, or prostate.
    \item \textbf{Vestibular schwannoma:} A benign tumour of the vestibular nerve, commonly located in the cerebellopontine angle, mimicking a meningioma in that region.
    \item \textbf{Solitary fibrous tumour / Haemangiopericytoma:} Highly vascular dural-based lesions that mimic meningioma radiologically but have a much higher tendency to recur and metastasise.
    \item \textbf{Dural lymphoma:} Primary or secondary CNS lymphoma presenting as a diffuse dural mass.
    \item \textbf{Inflammatory pachymeningitis:} Conditions such as neurosarcoidosis, tuberculosis, or IgG4-related disease, which cause localized thickening and inflammation of the dura mater.
\end{itemize}

\section*{When to Seek Medical Care}
Anyone experiencing new, persistent, or worsening headaches, progressive changes in vision, unexplained limb weakness, or difficulty with speech and balance should undergo a medical evaluation.

\textbf{Call 000 (or local emergency services) for:}
\begin{itemize}
    \item A new-onset seizure or sudden convulsive episode.
    \item Sudden, severe, or rapidly worsening headache (``thunderclap'' headache).
    \item Sudden weakness or numbness affecting one side of the body.
    \item Sudden loss of vision, speech difficulties, or loss of consciousness.
\end{itemize}

\section*{Diagnosis and Investigations}
The diagnosis of meningioma relies primarily on neuroimaging, followed by histopathological confirmation.
\begin{itemize}
    \item \textbf{Magnetic Resonance Imaging (MRI):} The diagnostic modality of choice. On MRI, meningiomas typically present as extra-axial, well-circumscribed masses that enhance homogeneously with gadolinium contrast. A characteristic ``dural tail sign'' (linear enhancement of the adjacent dura) is often visible.
    \item \textbf{Computed Tomography (CT) scan:} Useful for identifying calcification within the tumour and evaluating hyperostosis (bone thickening) or osteolysis of the adjacent skull, which is common in meningiomas.
    \item \textbf{Histopathological examination:} Definitive diagnosis and grading are obtained via tissue biopsy or surgical resection. Under the WHO classification:
        \begin{itemize}
            \item \textbf{Grade I (Benign):} Accounts for $\sim$80--85\% of cases, slow-growing with low recurrence rates.
            \item \textbf{Grade II (Atypical):} Accounts for $\sim$15--18\% of cases, characterized by increased mitotic activity and a higher risk of recurrence.
            \item \textbf{Grade III (Anaplastic/Malignant):} Accounts for $\sim$1--3\% of cases, highly aggressive with rapid growth and poor outcomes.
        \end{itemize}
\end{itemize}

\section*{Management}
The management strategy is highly individualised, depending on the patient's age, comorbidities, symptoms, and the size, location, and growth rate of the tumour.
\begin{itemize}
    \item \textbf{Active surveillance:} Suitable for patients with small, asymptomatic, incidentally discovered meningiomas. This involves serial brain MRIs (typically every 6 to 12 months initially) to monitor for changes.
    \item \textbf{Surgical resection:} The primary treatment modality for symptomatic or growing meningiomas. The goal is complete surgical removal of the tumour, its dural attachment, and any affected bone (Simpson Grade I resection), where anatomically safe.
    \item \textbf{Radiation therapy:} Standard external beam radiotherapy or stereotactic radiosurgery (SRS, such as Gamma Knife) is utilised for inoperable meningiomas, symptomatic residual tumours after subtotal resection, or as adjuvant therapy for Grade II (atypical) and Grade III (malignant) tumours.
    \item \textbf{Pharmacological therapy:} Corticosteroids (e.g., dexamethasone) are prescribed perioperatively to manage vasogenic oedema surrounding the tumour. Antiseizure medications (e.g., levetiracetam) are indicated if the patient has experienced seizures.
\end{itemize}

\section*{Complications}
Complications can arise from the mechanical pressure of the tumour itself, or as adverse effects of surgical and radiological treatments.
\begin{itemize}
    \item Vasogenic brain oedema, leading to increased intracranial pressure, focal deficits, or brain herniation.
    \text{Surgical complications:} Intracranial haemorrhage, surgical site infection, cerebrospinal fluid (CSF) leak, venous thromboembolism, or damage to adjacent cranial nerves.
    \item Radiation-induced complications: Radiation necrosis, localized alopecia, cognitive fatigue, or rare secondary malignancies.
    \item Permanent neurological deficits: Such as persistent weakness, speech difficulties, or visual deficits.
    \text{Epilepsy:} Persistent seizure disorders requiring long-term anticonvulsant therapy.
\end{itemize}

\section*{Prognosis and Outcomes}
The prognosis for patients with meningiomas is generally favourable, particularly for WHO Grade I tumours. Following a complete surgical resection, the 10-year survival rate exceeds 90\%, and the risk of recurrence is low (less than 10\% at 10 years). However, the prognosis declines for atypical (Grade II) and malignant (Grade III) meningiomas, which carry recurrence rates of up to 40\% and 80\% respectively, even after complete resection and adjuvant radiation therapy. Tumours located in anatomically challenging areas, such as the skull base or parasellar region, carry a higher risk of surgical morbidity and subtotal resection, which increases the likelihood of recurrence.

\section*{Prevention and Self-Care}
Because the exact cause of most meningiomas is unknown, there are few specific preventative measures.
\begin{itemize}
    \item Avoid unnecessary exposure to ionising radiation, particularly to the head.
    \item Inform healthcare providers of any history of meningioma before starting hormone replacement therapy or taking progesterone-containing medications.
    \item Adhere strictly to the recommended follow-up imaging schedule to monitor for recurrence or tumour growth.
    \item Participate in neurological rehabilitation (physiotherapy, occupational therapy, or speech pathology) if suffering from residual deficits post-treatment.
\end{itemize}

\section*{Sources and Further Reading}
The following reputable organisations provide further information and clinical guidance on meningiomas:
\begin{itemize}
    \item Cancer Council Australia. \textit{Information on brain cancer, meningiomas, and support services.} https://www.cancer.org.au/
    \item Brain Foundation (Australia). \textit{Fact sheets and research updates on brain tumours and meningiomas.} https://brainfoundation.org.au/
    \item NHS (UK). \textit{Patient information and clinical pathways for meningioma management.} https://www.nhs.uk/
    \item Mayo Clinic. \textit{Patient-focused overview of meningioma diagnosis, symptoms, and treatment.} https://www.mayoclinic.org/
    \item MSD Manual (Professional Edition). \textit{Clinical guide to intracranial tumours and meningiomas.} https://www.msdmanuals.com/
\end{itemize}
